% Section 7. Drafted from docs/PAPER-OUTLINE.md section 7 and docs/THEORY.md "Layer B". Every named
% result is a gated, axiom-free theorem in coq/*.v (TransportCSP.v and TransportCSPHard.v unless noted).

\section{Layer B: consistent states}\label{sec:layer-b}

Layer~B asks which states are consistent. It is a static question about a network's equations, not
about any run: when registries pass values along transport maps, a consistent assignment is one that
every map respects. This section shows that the question is a finite-domain constraint satisfaction
problem, that the fixed points of the repair dynamics are exactly its solutions, and that the
development's results sit at known points of the algebraic classification of such problems. Little of
the mathematics here is new; the section's role is to place consistency where a mature theory applies,
and to connect it to settlement through the interface of \cref{sec:interfaces}.

% ---------------------------------------------------------------------------------------------
\subsection{Consistency as constraint satisfaction}

Fix a finite family $F$ of transport maps on a finite set $D$. The \emph{template} $\Gamma_F$ consists
of the graphs $\{(x, f(x))\}$ of the maps in $F$, together with the unary relations $\{c\}$ for values
that some vertices are pinned to. A network whose edges are labeled by maps in $F$, with some vertices
pinned, is then an instance of $\mathrm{CSP}(\Gamma_F)$.

\begin{theorem}[\coq{section_iff_csp}, \coq{edge_fixed_iff_csp}]\label{thm:csp}
An assignment of values to the vertices of such a network is a section, a consistent assignment that
every edge's map respects, if and only if it is a solution of the corresponding instance of
$\mathrm{CSP}(\Gamma_F)$. Moreover, for a network without pinned vertices on a store with the usual
read and write laws, the states fixed by every edge repair of the network's dynamics are exactly those
solutions.
\end{theorem}

\emph{Proof idea.} The correspondence is unfolding: the constraint of an edge $u \to v$ labeled $f$ holds
at an assignment $s$ exactly when $f(s_u) = s_v$ (\coq{csat_econ}), and a pin constraint exactly when
the pinned value is taken (\coq{csat_pcon}), so sections agreeing with the pins are the solutions
(\coq{section_solution}); each constraint uses the graph of its edge's label or a pin, so the instance
is over $\Gamma_F$. Finiteness of $F$ and $D$ is not used. For the dynamics, the repair
$s[v := f(s_u)]$ fixes $s$ exactly when $f(s_u) = s_v$: reading $v$ after the write gives one
direction, and writing back the value just read leaves the store unchanged for the other
(\coq{erep_fixed_iff}). Taking every edge gives \coq{edge_fixed_iff_section}, and the correspondence turns
sections into solutions.

Merges in which a vertex reads several inputs are graphs of $k$-ary maps and fit the same way
(\coq{hsection_iff_csp}). The identification of global sections with constraint solutions is known in
general, from the sheaf-theoretic treatment of contextuality~\cite{abramsky2011,abramsky2014}; what is proved here
is its form for transport networks, together with the link to the dynamics.

An operation is a \emph{polymorphism} of the graph of $f$ exactly when it commutes with $f$
(\coq{pol_iff_commute}). So the polymorphisms of $\Gamma_F$ are the operations that commute with every
map in $F$ and fix every pinned value (\coq{pol_gamma_iff}), the centralizer of the family in the
sense of clone theory. By the dichotomy theorem for finite-domain constraint
problems~\cite{bulatov2017,zhuk2020}, cited and not mechanized, section existence for each fixed
finite family is either polynomial or NP-complete, and the polymorphisms decide which.

\emph{Proof idea.} Applied to $k$ rows $(x_i, f(x_i))$ of the graph of $f$, an operation $p$ that
preserves the graph returns a pair of the graph, which says $f(p(x_1, \dots, x_k)) = p(f(x_1), \dots,
f(x_k))$; conversely, commutation sends the column of left entries to a value whose image under $f$ is $p$ of
the column of right entries (\coq{pol_iff_commute}). For a pin $\{c\}$ the only rows are constant, so $p$
preserves it exactly when $p(c, \dots, c) = c$ (\coq{pol_pin_iff}), and \coq{pol_gamma_iff} is the
conjunction over the relations of the template.

% ---------------------------------------------------------------------------------------------
\subsection{Where the development's results sit}

\paragraph{Groups.} For group labels acting by translation, $x y^{-1} z$ is a Mal'tsev polymorphism
(\coq{group_maltsev}, \coq{group_pol}), the tractable case of Bulatov and Dalmau~\cite{bulatov2006};
the polynomial criterion mechanized for it is a coboundary test (\coq{section_iff_coboundary}).

\paragraph{Monotone maps on a chain: tractable.} A monotone map on a chain commutes with $\min$,
$\max$ and the median (\coq{chain_pol}). For pinned networks of monotone maps on a finite chain, arc
consistency decides whether a section exists, and the minimum of each arc-consistent domain is one
(\coq{ac_exact}). This is the binary, min-closed case of a classical result: max-closed
constraints~\cite{jeavons1995}, and more generally any semilattice polymorphism~\cite{jeavons1997},
make arc consistency decide the instance. Both hypotheses are needed. On the four-element diamond a
monotone map need not commute with join (\coq{diamond_join_fails}), and without monotonicity arc
consistency accepts an instance that has no section (\coq{ac_needs_mono}).

\emph{Proof idea.} The procedure revises every domain at once, keeping a value only if it has a
support along every edge into and out of its vertex, and stops when a pass removes nothing; every
other pass shrinks the total size of the domains, which starts at most the number of vertices times
$N + 1$, so it stops (\coq{measure_nxt}, \coq{ac_run_some}). Revision never removes the value of a
section (\coq{sol_ok}, \coq{run_spec}), so if a section exists every final domain is nonempty; this
direction does not use monotonicity. Conversely, take the minimum of each final domain
(\coq{ac_min_section}). For an edge $u \to v$ labeled $f$ with minima $a$ and $b$, the support of $a$
along the edge gives $f(a)$ in the domain of $v$, so $b \le f(a)$; the support of $b$ gives some $a'$
in the domain of $u$ with $f(a') = b$, and $a \le a'$ with monotonicity of $f$ gives $f(a) \le b$.
The pins hold because each domain starts inside its pinned value.

\paragraph{Monotone maps off the chain: hard.} Read as two flags, the diamond carries six monotone
maps (read one flag, their AND or OR, or a constant) whose only polymorphisms are projections, at
every arity (\coq{pol_dfam_iff}); so they have no Siggers, majority or Mal'tsev polymorphism
(\coq{dfam_no_siggers}). A linear reduction from 3-SAT to networks over these maps shows section
existence NP-complete (\coq{dnet_iff_sat}), for every family that contains them
(\coq{mono_family_hard}), and each of the six is needed (\coq{dfam_minimal}). This cell recovers known
boundaries: the six maps encode the two-element lattice, whose equation problem is
NP-complete~\cite{larose2006}, and for three-element ordered sets the boundary follows from the
classification of unary algebras on at most three elements~\cite{broniek2006}. What separates the two
cells is join preservation. Maps that preserve join, the shape of state-based CRDTs, keep join as a
semilattice polymorphism on any lattice (\coq{semilattice_pol}, \coq{joinhom_pol}); monotonicity alone
is not enough off the chain.

\emph{Proof idea for \coq{pol_dfam_iff}.} Projections are polymorphisms of every template of graphs
and pins (\coq{proj_pol}). Conversely, by \coq{pol_gamma_iff} a polymorphism $p$ commutes with the six
maps. Commuting with the maps that copy one flag into both shows that $p$ sends tuples from the chain
$\{(0,0), (1,1)\}$ into that chain, which defines a Boolean function $\sigma$ of $k$ arguments, and that
each flag of $p(x_1, \dots, x_k)$ is $\sigma$ applied to that flag of the arguments (\coq{p_btup},
\coq{sig_bit1}, \coq{sig_bit2}); commuting with the AND and OR maps makes $\sigma$ preserve $\wedge$
and $\vee$, and the two constants make it fix $0$ and $1$ (\coq{sig_comb}, \coq{sig_const}). Such a
$\sigma$ is a coordinate (\coq{bool_hom_proj}): writing each argument tuple $g$ as the join of its
unit vectors gives $\sigma(g) = \bigvee_j (g_j \wedge \sigma(e_j))$, so $\sigma(1, \dots, 1) = 1$
yields an $i$ with $\sigma(e_i) = 1$, and if $g_i = 0$ then $\sigma(g) = \sigma(g) \wedge \sigma(e_i)
= \sigma(g \wedge e_i) = \sigma(0, \dots, 0) = 0$. So both flags of $p$ are read from argument $i$,
and $p$ is the projection to it.

\paragraph{The reduction's own family.} The 3-SAT networks used for the lossy-network hardness result
use one fixed family of 26 maps, so constraint satisfaction over that single template is NP-complete
(\coq{hard_family_csp}).

\paragraph{Majority off the chain.} A family on pairs of Booleans that includes a lossy map commutes
with the coordinatewise median but not with meet or join (\coq{median_family}); tractability then
follows from the bounded-width results~\cite{jeavons1997,federvardi1998}, cited, not from a mechanized
algorithm.

% ---------------------------------------------------------------------------------------------
\subsection{What is open}

Whether every finite ordered set with at least three elements that is not a chain has only
projections as polymorphisms of its monotone maps is open in general as far as the literature checked
shows. The bound on size is needed: on the two-element antichain every self-map is monotone, and the
majority and minority operations commute with all of them. The diamond is proved here;
an argument for bounded ordered sets is recorded in the development but not mechanized; three- and
four-element cases can be read off published classifications of centralizing monoids; and ordered
sets without a top and a bottom, from five elements up, are not settled by any source found.
